%%%%%%%%%%%%%%%%
%%% exod 18 %%%
%%%%%%%%%%%%%%%%

\Chapter{18}

\Verse{1}
{
	\hE{וַיִּשְׁמַע יִתְרוֹ כֹהֵן מִדְיָן חֹתֵן מֹשֶׁה אֵת כּל אֲשֶׁר עָשָׂה אֱלֹהִים לְמֹשֶׁה וּלְיִשְׂרָאֵל עַמּוֹ כִּי הוֹצִיא יְהֹוָה אֶת יִשְׂרָאֵל מִמִּצְרָיִם׃}
}
{
	\eE{
		And Jethro, priest of Midian, Moses’ father-in-law, heard
		everything that {\God} had done for Moses and for Israel, His
		people, that YHWH had brought Israel out from Egypt.
	}
}
{
}

\Verse{2}
{
	\hE{וַיִּקַּח יִתְרוֹ חֹתֵן מֹשֶׁה אֶת צִפֹּרָה אֵשֶׁת מֹשֶׁה}
	\hRJE{אַחַר שִׁלּוּחֶיהָ׃}
}
{
	\eE{
		And Jethro, Moses’ father-in-law, took Zipporah, Moses’
		wife,}
	\eRJE{after her being sent off,}
	}
}
{
}

\Verse{3}
{
	\hE{וְאֵת שְׁנֵי בָנֶיהָ אֲשֶׁר שֵׁם הָאֶחָד גֵּרְשֹׁם כִּי אָמַר גֵּר הָיִיתִי בְּאֶרֶץ נכְרִיָּה׃}
}
{
	\eE{
		and her two sons, of whom one’s name was Gershom, because
		he said, “I was an alien in a foreign land,”
	}
}
{
}

\Verse{4}
{
	\hE{וְשֵׁם הָאֶחָד אֱלִיעֶזֶר כִּי אֱלֹהֵי אָבִי בְּעֶזְרִי וַיַּצִּלֵנִי מֵחֶרֶב פַּרְעֹה׃}
}
{
	\eE{
		and one’s name was Eliezer, “because my father’s {\God} was
		my help and rescued me from Pharaoh’s sword.”
	}
}
{
}

\Verse{5}
{
	\hE{וַיָּבֹא יִתְרוֹ חֹתֵן מֹשֶׁה וּבָנָיו וְאִשְׁתּוֹ אֶל מֹשֶׁה אֶל הַמִּדְבָּר אֲשֶׁר הוּא חֹנֶה שָׁם הַר הָאֱלֹהִים׃}
}
{
	\eE{
		And Jethro, Moses’ father-in-law, and his sons and his
		wife came to Moses, to the wilderness in which he was
		camping, at the Mountain of {\God}.
	}
}
{
}

\Verse{6}
{
	\hE{וַיֹּאמֶר אֶל מֹשֶׁה אֲנִי חֹתֶנְךָ יִתְרוֹ בָּא אֵלֶיךָ וְאִשְׁתְּךָ וּשְׁנֵי בָנֶיהָ עִמָּהּ׃}
}
{
	\eE{
		And he said to Moses, “I, your father-in-law, Jethro, have
		come to you, and your wife and her two sons with her.”
	}
}
{
}

\Verse{7}
{
	\hE{וַיֵּצֵא מֹשֶׁה לִקְרַאת חֹתְנוֹ וַיִּשְׁתַּחוּ וַיִּשַּׁק לוֹ וַיִּשְׁאֲלוּ אִישׁ לְרֵעֵהוּ לְשָׁלוֹם וַיָּבֹאוּ הָאֹהֱלָה׃}
}
{
	\eE{
		And Moses went out to his father-in-law, and he bowed, and
		he kissed him, and they asked each other how they were,
		and they came to the tent.
	}
}
{
%	\footnote{
%		Moses went out to his father-in-law. Jethro arrives with Moses’ wife and
%		sons, but only Moses’ meeting with Jethro is described. There is nothing
%		about his reunion with his family. The Bible’s silences are frustrating
%		(where did Cain find his wife? how did Sarah feel about her son’s being
%		sacrificed?), but I think that we should learn to live with them. Often
%		they are very loud silences—which is to say: very powerful. Why does the
%		Torah announce the arrival of Moses’ wife and children but not describe
%		a personal meeting scene between them and Moses? Presumably, this is
%		because it was important to let us know that his wife and children were
%		present with the Israelites during the wilderness period, but it was not
%		deemed important to describe their family relations. Family relations
%		are described in the age of the patriarchs and in the age of the kings
%		because, there, they play a role in the destiny of the people. They are
%		described in the cases of Jephthah and Samson because they propel the
%		life events of those persons. But in Moses’ case the family episodes are
%		few: his mother and sister protect him as a baby, his wife has a
%		confrontation with him at the inn on the way to Egypt, Aaron and Miriam
%		criticize his Cushite wife and are rebuked by God. The rest of Moses’
%		family life is personal. His wife speaks three sentences in the Torah.
%		His sons speak none. That tells us, at the very least, that his
%		particular life mission is primarily about him and God, on the one hand,
%		and him and the people, on the other. And the nature of that mission is
%		such that his family situation is not permitted to impinge upon it. The
%		missions of Abraham, David, and the rest of us are not like that. It is
%		our frustration with the silences and our hunger for more Torah that
%		give birth to midrash. So I was once asked to make a midrash on what
%		might have happened between Moses and his wife and sons. I suggested: He
%		told them everything—the confrontations with Pharaoh, the snakes, the
%		plagues, the splitting of the sea. And his sons asked him why he had
%		sent them away and thus kept them from seeing all that. And he told them
%		that it was dangerous, and he was protecting them. Did that show
%		insufficient faith on his part that things would go well in Egypt? Of
%		course, he says. But he was just starting out and did not yet know the
%		divine power—or the divine character. And so a father’s love was his
%		first motive. But now, he assured them, they would stay with him and get
%		to see an even greater event: Sinai.
%	}
}

\Verse{8}
{
	\hE{וַיְסַפֵּר מֹשֶׁה לְחֹתְנוֹ אֵת כּל אֲשֶׁר עָשָׂה יְהֹוָה לְפַרְעֹה וּלְמִצְרַיִם עַל אוֹדֹת יִשְׂרָאֵל אֵת כּל הַתְּלָאָה אֲשֶׁר מְצָאָתַם בַּדֶּרֶךְ וַיַּצִּלֵם יְהֹוָה׃}
}
{
	\eE{
		And Moses told his father-in-law everything that YHWH had
		done to Pharaoh and to Egypt with regard to Israel, all
		the hardship that had found them on the way, and YHWH
		rescued them.
	}
}
{
}

\Verse{9}
{
	\hE{וַיִּחַדְּ יִתְרוֹ עַל כּל הַטּוֹבָה אֲשֶׁר עָשָׂה יְהֹוָה לְיִשְׂרָאֵל אֲשֶׁר הִצִּילוֹ מִיַּד מִצְרָיִם׃}
}
{
	\eE{
		And Jethro rejoiced over all the good that YHWH had done
		for Israel, that He had rescued it from Egypt’s hand.
	}
}
{
}

\Verse{10}
{
	\hE{וַיֹּאמֶר יִתְרוֹ בָּרוּךְ יְהֹוָה אֲשֶׁר הִצִּיל אֶתְכֶם מִיַּד מִצְרַיִם וּמִיַּד פַּרְעֹה אֲשֶׁר הִצִּיל אֶת הָעָם מִתַּחַת יַד מִצְרָיִם׃}
}
{
	\eE{
		And Jethro said, “Blessed is YHWH, who rescued you from
		Egypt’s hand and from Pharaoh’s hand, who rescued the
		people from under Egypt’s hand:
	}
}
{
}

\Verse{11}
{
	\hE{עַתָּה יָדַעְתִּי כִּי גָדוֹל יְהֹוָה מִכּל הָאֱלֹהִים כִּי בַדָּבָר אֲשֶׁר זָדוּ עֲלֵיהֶם׃}
}
{
	\eE{
		now I know that YHWH is bigger than all the gods, because
		of the thing that they plotted against them.”
	}
}
{
%	\footnote{
%		the thing that they plotted against them. Meaning: Egypt’s plot to
%		control the Israelites by enslaving them.
%	}
}

\Verse{12}
{
	\hE{וַיִּקַּח יִתְרוֹ חֹתֵן מֹשֶׁה עֹלָה וּזְבָחִים לֵאלֹהִים וַיָּבֹא אַהֲרֹן וְכֹל זִקְנֵי יִשְׂרָאֵל לֶאֱכל לֶחֶם עִם חֹתֵן מֹשֶׁה לִפְנֵי הָאֱלֹהִים׃}
}
{
	\eE{
		And Jethro, Moses’ father-in-law, took a burnt offering
		and sacrifices to {\God}. And Aaron and all of Israel’s
		elders came to eat bread with Moses’ father-in-law before
		{\God}.
	}
}
{
}

\Verse{13}
{
	\hE{וַיְהִי מִמּחֳרָת וַיֵּשֶׁב מֹשֶׁה לִשְׁפֹּט אֶת הָעָם וַיַּעֲמֹד הָעָם עַל מֹשֶׁה מִן הַבֹּקֶר עַד הָעָרֶב׃}
}
{
	\eE{
		And it was the next day, and Moses sat to judge the
		people. And the people stood by Moses from the morning to
		the evening.
	}
}
{
}

\Verse{14}
{
	\hE{וַיַּרְא חֹתֵן מֹשֶׁה אֵת כּל אֲשֶׁר הוּא עֹשֶׂה לָעָם וַיֹּאמֶר מָה הַדָּבָר הַזֶּה אֲשֶׁר אַתָּה עֹשֶׂה לָעָם מַדּוּעַ אַתָּה יוֹשֵׁב לְבַדֶּךָ וְכל הָעָם נִצָּב עָלֶיךָ מִן בֹּקֶר עַד עָרֶב׃}
}
{
	\eE{
		And Moses’ father-in-law saw all that he was doing for the
		people, and he said, “What’s this thing that you’re doing
		for the people? Why are you sitting by yourself, and the
		entire people is standing up by you from morning until
		evening?”
	}
}
{
}

\Verse{15}
{
	\hE{וַיֹּאמֶר מֹשֶׁה לְחֹתְנוֹ כִּי יָבֹא אֵלַי הָעָם לִדְרֹשׁ אֱלֹהִים׃}
}
{
	\eE{
		And Moses said to his father-in-law, “Because the people
		come to me to inquire of {\God}.
	}
}
{
}

\Verse{16}
{
	\hE{כִּי יִהְיֶה לָהֶם דָּבָר בָּא אֵלַי וְשָׁפַטְתִּי בֵּין אִישׁ וּבֵין רֵעֵהוּ וְהוֹדַעְתִּי אֶת חֻקֵּי הָאֱלֹהִים וְאֶת תּוֹרֹתָיו׃}
}
{
	\eE{
		When they have a matter, it comes to me, and I judge
		between each one and his companion, and I make known {\God}’s
		laws and His instructions.”
	}
}
{
}

\Verse{17}
{
	\hE{וַיֹּאמֶר חֹתֵן מֹשֶׁה אֵלָיו לֹא טוֹב הַדָּבָר אֲשֶׁר אַתָּה עֹשֶׂה׃}
}
{
	\eE{
		And Moses’ father-in-law said to him, “The thing that
		you’re doing isn’t good.
	}
}
{
}

\Verse{18}
{
	\hE{נָבֹל תִּבֹּל גַּם אַתָּה גַּם הָעָם הַזֶּה אֲשֶׁר עִמָּךְ כִּי כָבֵד מִמְּךָ הַדָּבָר לֹא תוּכַל עֲשֹׂהוּ לְבַדֶּךָ׃}
}
{
	\eE{
		You’ll be worn out, both you and this people who are with
		you, because the thing is too heavy for you. You won’t be
		able to do it by yourself.
	}
}
{
%	\footnote{
%		too heavy for you. “Heavy” again. See the comment on Exod 5:9. And see
%		Num 11:14. Footnote 18:18. able to do it by yourself. The convergence of
%		the terms for “to be by oneself” and “to be able” (twice: vv. 18 and 23)
%		in this passage is striking because they both appear in the story of
%		Jacob wrestling with God (Gen 32:24–33), which begins with the notice
%		that “he was left by himself” and later notes (twice) that Jacob “was
%		able” in the struggle. It is remarkable that Jacob was pictured as able
%		to struggle with God by himself but Moses is pictured as not being able
%		to manage the people by himself. Does a person stand a better chance of
%		being successful in a personal struggle with God than in taking on a
%		community of humans?
%	}
}

\Verse{19}
{
	\hE{עַתָּה שְׁמַע בְּקֹלִי אִיעָצְךָ וִיהִי אֱלֹהִים עִמָּךְ הֱיֵה אַתָּה לָעָם מוּל הָאֱלֹהִים וְהֵבֵאתָ אַתָּה אֶת הַדְּבָרִים אֶל הָאֱלֹהִים׃}
}
{
	\eE{
		Now listen to my voice, and I’ll advise you, and may {\God}
		be with you. You be for the people toward {\God}, and you
		will bring the matters to {\God}.
	}
}
{
}

\Verse{20}
{
	\hE{וְהִזְהַרְתָּה אֶתְהֶם אֶת הַחֻקִּים וְאֶת הַתּוֹרֹת וְהוֹדַעְתָּ לָהֶם אֶת הַדֶּרֶךְ יֵלְכוּ בָהּ וְאֶת הַמַּעֲשֶׂה אֲשֶׁר יַעֲשׂוּן׃}
}
{
	\eE{
		And you’ll enlighten them with the laws and the
		instructions, and you’ll make known to them the way in
		which they’ll go and the thing that they’ll do.
	}
}
{
}

\Verse{21}
{
	\hE{וְאַתָּה תֶחֱזֶה מִכּל הָעָם אַנְשֵׁי חַיִל יִרְאֵי אֱלֹהִים אַנְשֵׁי אֱמֶת שֹׂנְאֵי בָצַע וְשַׂמְתָּ עֲלֵהֶם שָׂרֵי אֲלָפִים שָׂרֵי מֵאוֹת שָׂרֵי חֲמִשִּׁים וְשָׂרֵי עֲשָׂרֹת׃}
}
{
	\eE{
		And you will envision, out of all the people, worthy men,
		who fear {\God}, men of truth, who hate bribery, and you’ll
		set chiefs of thousands, chiefs of hundreds, chiefs of
		fifties, and chiefs of tens over them.
	}
}
{
}

\Verse{22}
{
	\hE{וְשָׁפְטוּ אֶת הָעָם בְּכל עֵת וְהָיָה כּל הַדָּבָר הַגָּדֹל יָבִיאוּ אֵלֶיךָ וְכל הַדָּבָר הַקָּטֹן יִשְׁפְּטוּ הֵם וְהָקֵל מֵעָלֶיךָ וְנָשְׂאוּ אִתָּךְ׃}
}
{
	\eE{
		And they’ll judge the people at all times. And it will be:
		they’ll bring every matter that is big to you, and they
		will judge every matter that is small. And make it lighter
		on you, and they’ll bear it with you.
	}
}
{
}

\Verse{23}
{
	\hE{אִם אֶת הַדָּבָר הַזֶּה תַּעֲשֶׂה וְצִוְּךָ אֱלֹהִים וְיָכלְתָּ עֲמֹד וְגַם כּל הָעָם הַזֶּה עַל מְקֹמוֹ יָבֹא בְשָׁלוֹם׃}
}
{
	\eE{
		If you’ll do this thing, and YHWH will command you, then
		you’ll be able to stand, and also this entire people will
		come to its place in peace.”
	}
}
{
}

\Verse{24}
{
	\hE{וַיִּשְׁמַע מֹשֶׁה לְקוֹל חֹתְנוֹ וַיַּעַשׂ כֹּל אֲשֶׁר אָמָר׃}
}
{
	\eE{
		And Moses listened to his father-in-law’s voice and did
		everything that he had said.
	}
}
{
}

\Verse{25}
{
	\hE{וַיִּבְחַר מֹשֶׁה אַנְשֵׁי חַיִל מִכּל יִשְׂרָאֵל וַיִּתֵּן אֹתָם רָאשִׁים עַל הָעָם שָׂרֵי אֲלָפִים שָׂרֵי מֵאוֹת שָׂרֵי חֲמִשִּׁים וְשָׂרֵי עֲשָׂרֹת׃}
}
{
	\eE{
		And Moses chose worthy men out of all of Israel and made
		them heads over the people: chiefs of thousands, chiefs of
		hundreds, chiefs of fifties, and chiefs of tens.
	}
}
{
}

\Verse{26}
{
	\hE{וְשָׁפְטוּ אֶת הָעָם בְּכל עֵת אֶת הַדָּבָר הַקָּשֶׁה יְבִיאוּן אֶל מֹשֶׁה וְכל הַדָּבָר הַקָּטֹן יִשְׁפּוּטוּ הֵם׃}
}
{
	\eE{
		And they judged the people at all times. They would bring
		the matter that was hard to Moses, and they would judge
		every matter that was small.
	}
}
{
}

\Verse{27}
{
	\hE{וַיְשַׁלַּח מֹשֶׁה אֶת חֹתְנוֹ וַיֵּלֶךְ לוֹ אֶל אַרְצוֹ׃}
}
{
	\eE{
		And Moses let his father-in-law go, and he went to his
		land.
	}
}
{
}
